\section{Preliminaries}\label{sec:prelim}

\paragraph{Notation.}
Let $\Xcal$ be an arbitrary input space and let
$\Hcal\subset\{0,1\}^{\Xcal}$ be a hypothesis class of binary classifiers.
Examples are drawn from $\Zcal:=\Xcal\times\{0,1\}$ according to an unknown
distribution $\Dcal$. We write $S=\{(x_i,y_i)\}_{i=1}^{n}\sim\Dcal^{n}$ for an
i.i.d.\ sample. For a measurable $g:\Zcal\to[0,1]$ we use $\Dcal g:=\E_{(x,y)
\sim\Dcal}[g(x,y)]$ for its population mean and $\widehat\Dcal_n g:=\frac1n
\sum_{i=1}^{n} g(x_i,y_i)$ for its empirical mean on $S$.
Throughout, $c,C,C_1,C_2,\dots$ denote universal positive constants whose
values may change from line to line and depend on no problem parameter. We adopt
the standard measurability convention of empirical-process theory: $\Hcal$ is
assumed \emph{image-admissible Suslin} (equivalently, pointwise-separable), so
that suprema such as $\sup_{h\in\Hcal}\abs{\Risk(h)-\Riskhat_n(h)}$ are
measurable functions of $S$ and the concentration and symmetrization arguments
below apply verbatim.

\begin{definition}[Population and empirical risk]\label{def:risk}
For $h\in\Hcal$ the population and empirical (training) risks under the $0/1$
loss are
\begin{align}\label{eq:risk-def}
\Risk(h)
:=\Pr_{(x,y)\sim\Dcal}\bigl[h(x)\neq y\bigr]
=\E_{(x,y)\sim\Dcal}\bigl[\1\{h(x)\neq y\}\bigr],
\qquad
\Riskhat_n(h):=\frac1n\sum_{i=1}^{n}\1\{h(x_i)\neq y_i\}.
\end{align}
\end{definition}

\begin{definition}[$0/1$ loss class]\label{def:loss-class}
The loss class associated with $\Hcal$ is
\begin{align}\label{eq:loss-class}
\Lcal:=\bigl\{\,\ell_h:(x,y)\mapsto\1\{h(x)\neq y\}\ \big|\ h\in\Hcal\,\bigr\}
\subset\{0,1\}^{\Zcal}.
\end{align}
With this notation $\Risk(h)=\Dcal\ell_h$ and $\Riskhat_n(h)=\widehat\Dcal_n
\ell_h$, so that
$\sup_{h\in\Hcal}\abs{\Risk(h)-\Riskhat_n(h)}
=\sup_{\ell\in\Lcal}\abs{\Dcal\ell-\widehat\Dcal_n\ell}$.
\end{definition}

\begin{definition}[Empirical Rademacher complexity]\label{def:rademacher}
Let $\sgn_1,\dots,\sgn_m$ be i.i.d.\ Rademacher variables, i.e.\
$\Pr[\sgn_i=+1]=\Pr[\sgn_i=-1]=\tfrac12$, independent of the data. For a class
$\Gcal\subset\R^{\Zcal}$ and a sample $T=(z_1,\dots,z_m)$, the empirical and
average Rademacher complexities are
\begin{align}\label{eq:rad-def}
\Radhat_T(\Gcal):=\E_{\sgn}\Bigl[\sup_{g\in\Gcal}\frac1m\sum_{i=1}^{m}
   \sgn_i\,g(z_i)\Bigr],
\qquad
\Rad_m(\Gcal):=\E_{T\sim\Dcal^{m}}\bigl[\Radhat_T(\Gcal)\bigr].
\end{align}
\end{definition}

\begin{definition}[Growth function]\label{def:growth}
The growth function (shatter coefficient) of $\Hcal$ at size $m$ is
\begin{align}\label{eq:growth-def}
\Growth_{\Hcal}(m):=\max_{x_1,\dots,x_m\in\Xcal}
   \bigl|\bigl\{(h(x_1),\dots,h(x_m)):h\in\Hcal\bigr\}\bigr|.
\end{align}
A set $\{x_1,\dots,x_m\}$ is \emph{shattered} by $\Hcal$ if the maximum above
equals $2^{m}$, i.e.\ all $2^{m}$ label patterns are realized. The
Vapnik--Chervonenkis dimension $\VC(\Hcal)$ is the largest $m$ for which some
size-$m$ set is shattered.
\end{definition}

\begin{assumption}[I.i.d.\ sampling]\label{ass:iid}
The sample $S=\{(x_i,y_i)\}_{i=1}^{n}$ consists of $n$ independent draws from a
single distribution $\Dcal$ on $\Zcal$.
\end{assumption}

\begin{assumption}[Finite VC dimension, non-degenerate regime]\label{ass:vc}
The class $\Hcal$ has $\VC(\Hcal)=d<\infty$, and the sample size satisfies
$n\ge d\ge 1$.
\end{assumption}

\begin{remark}\label{rem:regime}
The hypothesis $n\ge d$ in \Cref{ass:vc} costs nothing. Since both $\Risk(h)$ and
$\Riskhat_n(h)$ lie in $[0,1]$, the deterministic trivial bound
$\sup_{h}\abs{\Risk(h)-\Riskhat_n(h)}\le 1$ holds for every sample and every $n$;
in the excluded regime $n<d$ one simply invokes this trivial bound, so no claim
about \Cref{thm:vc-bound} is needed there. The regime $n\ge d$ is exactly where
the Sauer--Shelah estimate $\Growth_{\Hcal}(m)\le(em/d)^{d}$ is informative and
the bound of \Cref{thm:vc-bound} is the operative one.
\end{remark}

\begin{fact}[The loss class has the same VC dimension]\label{fac:loss-vc}
$\Growth_{\Lcal}(m)\le\Growth_{\Hcal}(m)$ for all $m$, and hence
$\VC(\Lcal)\le\VC(\Hcal)=d$.
\end{fact}

\begin{proof}
Fix $z_1=(x_1,y_1),\dots,z_m=(x_m,y_m)\in\Zcal$. The map
\begin{align*}
\psi:(h(x_1),\dots,h(x_m))\ \longmapsto\
\bigl(\1\{h(x_1)\neq y_1\},\dots,\1\{h(x_m)\neq y_m\}\bigr),
\end{align*}
is a coordinatewise bijection of $\{0,1\}^{m}$: on coordinate $i$ it is the
identity when $y_i=0$ and the bit-flip $b\mapsto 1-b$ when $y_i=1$, both
bijections of $\{0,1\}$. A bijection cannot increase the number of distinct
images, so the label patterns of $\Lcal$ on $(z_1,\dots,z_m)$ are the image of
those of $\Hcal$ on $(x_1,\dots,x_m)$ under this bijection, giving
$\abs{\Lcal_{|(z_1,\dots,z_m)}}=\abs{\Hcal_{|(x_1,\dots,x_m)}}
\le\Growth_{\Hcal}(m)$. Taking the maximum over $z_1,\dots,z_m$ yields
$\Growth_{\Lcal}(m)\le\Growth_{\Hcal}(m)$, and a shattered set of $\Lcal$ would
force one of $\Hcal$ of the same size, so $\VC(\Lcal)\le d$. This gives the
assertion.
\end{proof}
